%% toc: P1 1 1.3,1.4
\section{Formula schemas, metavariables, quotation, and codes}
\label{sec:p01-w02-p01-004-schemas-codes}

\Cref{sec:levels-02} separates the object language from the metalanguage and
notes that a Greek letter such as \(\varphi\) is normally a metavariable,
not an object-language symbol. That remark leaves three devices undeveloped:
the formula schema, in which metavariables stand for whole families of
formulas; the quotation name, which mentions one inscription; and the
numerical code, which assigns one natural number to one formula and, in a
language with numerals, lets the object language itself name that formula.
This section develops all three for the propositional language and states
exactly how they differ.

\paragraph{Inputs.} The propositional alphabet, the formula set
\(\mathsf{Fm}_{\mathsf{CPL}}\) defined as the least set of strings closed
under the formation clauses for \(p_n\), \(\bottom\), \(\neg\) and the
parenthesized binary connectives \(\land,\lor,\to\), and the unique
readability theorem \cref{thm:prop-unique-readability}
\cite[ch.~``Propositional Syntax and Semantics'', definitions
``Propositional alphabet'' and ``Propositional formula'', theorem ``Unique
readability'']{FrozenFormalLogicSchemasCodes}; the numerical pairing
\(\operatorname{pair}_{\N}(a,b)=\frac{(a+b)(a+b+1)}{2}+b\) and its proof of
bijectivity \(\N\times\N\to\N\) in \cref{sec:strings-01}
\cite[ch.~``Finite Strings, Induction, and Recursion'', section ``Finite
sequences and strings'']{FrozenFormalLogicSchemasCodes}; and, for the
relationships at the end, uniform substitution
(\cref{def:p02-uniform-substitution}), the arithmetic coding of
\cref{sec:arithmetic-coding-04} and \cref{thm:tarski-undefinability}.

\paragraph{Conclusions.} A definition of formula schemas and their instances
(\cref{def:p01-w02-p01-004-schema,def:p01-w02-p01-004-instance}) with unique
readability for schemas (\cref{thm:p01-w02-p01-004-schema-readability}); the
uniqueness of matching (\cref{thm:p01-w02-p01-004-matching-unique}), a
matching procedure and its correctness
(\cref{def:p01-w02-p01-004-matching,thm:p01-w02-p01-004-matching-correct});
the size of instance sets (\cref{thm:p01-w02-p01-004-instance-count}) and their
relation to uniform substitution
(\cref{thm:p01-w02-p01-004-substitution-link}); a numerical code for
formulas with injectivity and decrease on immediate subformulas
(\cref{def:p01-w02-p01-004-code,thm:p01-w02-p01-004-code-properties}); a
decoding procedure deciding which numbers are codes
(\cref{def:p01-w02-p01-004-decoding,thm:p01-w02-p01-004-decoding-correct});
and the schema/code distinction stated as a theorem
(\cref{def:p01-w02-p01-004-distinction,thm:p01-w02-p01-004-distinction}),
with a decidable code set for each instance family
(\cref{thm:p01-w02-p01-004-instance-codes}). All arguments are editorial
completions built on the inherited definitions; the proofs are in the annex.

\subsection{Metavariables and formula schemas}
\label{sec:p01-w02-p01-004-schemas}

A metavariable such as \(\varphi\) is a variable of the metalanguage whose
values are formulas. In a sentence such as ``for all formulas \(\varphi\),
the formula \((\varphi\to\varphi)\) is valid'' it is bound by a metalinguistic
quantifier and takes one formula as its value at a time. A display such as
\((\varphi\to(\psi\to\varphi))\), offered as an axiom or as the premise
pattern of a rule, is used differently: it is read as a pattern, and the
letters in it mark the places where formulas are to be inserted. To make the
pattern itself an object of study, the letters are taken as additional
symbols of a metalinguistic syntax.

\begin{hbdefinition}{def:p01-w02-p01-004-schema}{Schematic letters and formula schemas}
Fix a countably infinite set \(\Xi=\{\xi_k:k\in\N\}\) of \emph{schematic
letters}, pairwise distinct and distinct from every symbol of the
propositional alphabet. Write \(A\) for the propositional alphabet and
\(A^{+}=A\cup\Xi\). The set \(\mathrm{Sch}\) of \emph{formula schemas} is the
least set of strings over \(A^{+}\) such that:
\begin{enumerate}
  \item \(\xi_k\in\mathrm{Sch}\) for every \(k\in\N\);
  \item \(p_n\in\mathrm{Sch}\) for every \(n\in\N\), and
        \(\bottom\in\mathrm{Sch}\);
  \item if \(\mathfrak s\in\mathrm{Sch}\), then \(\neg\mathfrak s\in\mathrm{Sch}\);
  \item if \(\mathfrak s,\mathfrak t\in\mathrm{Sch}\) and
        \(\circ\in\{\land,\lor,\to\}\), then
        \((\mathfrak s\circ\mathfrak t)\in\mathrm{Sch}\).
\end{enumerate}
For a schema \(\mathfrak s\), \(\operatorname{MV}(\mathfrak s)\) is the set of
schematic letters occurring in \(\mathfrak s\); it is finite because
\(\mathfrak s\) is a finite string. The letters \(\xi_k\) are symbols of the
metalanguage. They are not added to the object language, and a schema
containing one of them is not a formula.
\end{hbdefinition}

In displays, the handbook writes schemas with the formula metavariables
\(\varphi,\psi,\chi\) in place of \(\xi_0,\xi_1,\xi_2\), and with the
typographical omissions of outer parentheses declared for formulas in
\cref{ch:prop-syntax-semantics}. Thus the display
\(\varphi\to(\psi\to\varphi)\), read as a schema, abbreviates the string
\((\xi_0\to(\xi_1\to\xi_0))\).

\begin{hbtheorem}{thm:p01-w02-p01-004-schema-readability}{proved=annex:p01-w02-p01-004-schema-readability}{Unique readability for schemas}
Every member of \(\mathrm{Sch}\) satisfies exactly one of the four formation
clauses, and in the first two clauses it is a single token. In the negation
case its immediate subschema is unique. In a binary case its main connective
and its two immediate subschemas are unique. Moreover
\(\mathsf{Fm}_{\mathsf{CPL}}=\{\mathfrak s\in\mathrm{Sch}:
\operatorname{MV}(\mathfrak s)=\varnothing\}\).
\end{hbtheorem}

\begin{hbdefinition}{def:p01-w02-p01-004-instance}{Instantiation and instance}
An \emph{instantiation} is a function
\(\iota:\Xi\to\mathsf{Fm}_{\mathsf{CPL}}\). Its action on schemas is defined
by structural recursion, which is well defined by
\cref{thm:p01-w02-p01-004-schema-readability}:
\[
\begin{aligned}
 \xi_k[\iota]&=\iota(\xi_k),&
 p_n[\iota]&=p_n,&
 \bottom[\iota]&=\bottom,\\
 (\neg\mathfrak s)[\iota]&=\neg(\mathfrak s[\iota]),&
 (\mathfrak s\circ\mathfrak t)[\iota]
   &=(\mathfrak s[\iota]\circ\mathfrak t[\iota]).
\end{aligned}
\]
The formula \(\mathfrak s[\iota]\) is the \emph{instance of \(\mathfrak s\)
under \(\iota\)}. The \emph{instance set} of \(\mathfrak s\) is
\(\operatorname{Inst}(\mathfrak s)=\{\mathfrak s[\iota]:\iota
\text{ an instantiation}\}\). A formula \(\varphi\) \emph{matches}
\(\mathfrak s\) when \(\varphi\in\operatorname{Inst}(\mathfrak s)\).
\end{hbdefinition}

Induction on schemas shows that \(\mathfrak s[\iota]\in
\mathsf{Fm}_{\mathsf{CPL}}\) for every \(\iota\), because each clause applies a
formation clause of \(\mathsf{Fm}_{\mathsf{CPL}}\) to formulas. Every
occurrence of the same letter receives the same formula: instantiation is
uniform. An axiom schema of a calculus contributes all members of its
instance set as axioms, and a rule schema licenses all the rule instances
obtained by instantiating its premises and conclusion with one common
\(\iota\); a side condition attached to a rule further restricts which
instantiations the rule licenses.

\begin{hbconvention}{def:p01-w02-p01-004-display-reading}{The two readings of a display containing metavariables}
Let \(\varphi_0,\ldots,\varphi_{m}\) be metavariables with current values, as
inside the scope of a metalinguistic quantifier ``for all formulas
\(\varphi_0,\ldots,\varphi_m\)''. A display built from them by the formation
clauses denotes the single formula \(\mathfrak s[\iota]\), where
\(\mathfrak s\) is the schema obtained by writing \(\xi_j\) for
\(\varphi_j\) and \(\iota(\xi_j)\) is the current value of \(\varphi_j\); this
is the \emph{bound reading}. A display offered as an axiom, a rule premise or
a rule conclusion, with no enclosing quantifier, denotes its schema, and the
claim made with it concerns every member of \(\operatorname{Inst}(\mathfrak
s)\); this is the \emph{schematic reading}. A claim stated in the schematic
reading is equivalent to the claim prefixed by ``for all formulas
\(\varphi_0,\ldots,\varphi_m\)'' in the bound reading, by the definition of
\(\operatorname{Inst}\).
\end{hbconvention}

\begin{hbtheorem}{thm:p01-w02-p01-004-matching-unique}{proved=annex:p01-w02-p01-004-matching-unique}{Uniqueness of matching}
Let \(\mathfrak s\in\mathrm{Sch}\) and let \(\iota,\iota'\) be
instantiations.
\begin{enumerate}
  \item If \(\iota\) and \(\iota'\) agree on \(\operatorname{MV}(\mathfrak
        s)\), then \(\mathfrak s[\iota]=\mathfrak s[\iota']\).
  \item If \(\mathfrak s[\iota]=\mathfrak s[\iota']\), then \(\iota\) and
        \(\iota'\) agree on \(\operatorname{MV}(\mathfrak s)\).
\end{enumerate}
Hence a formula that matches \(\mathfrak s\) determines the values of the
instantiation on the letters of \(\mathfrak s\) uniquely.
\end{hbtheorem}

The second clause depends on the unique readability of formulas. Without the
parentheses required by the binary formation clause, the string
\(p_0\land p_1\land p_2\) would match the pattern
\(\varphi\land\psi\) in two ways, with \(\varphi\) taking the value
\(p_0\) or the value \(p_0\land p_1\), and the instantiation would not be
recoverable from the instance.

\begin{hbmethod}{def:p01-w02-p01-004-matching}{Matching a formula against a schema}
\emph{Input:} a schema \(\mathfrak s\in\mathrm{Sch}\) and a formula
\(\varphi\in\mathsf{Fm}_{\mathsf{CPL}}\). \emph{Output:} a finite function
\(m\) from schematic letters to formulas, or the answer \(\mathsf{fail}\).
The value \(\operatorname{match}(\mathfrak s,\varphi)\) is defined by
recursion on \(\mathfrak s\):
\begin{enumerate}
  \item If \(\mathfrak s=\xi_k\), return \(\{\xi_k\mapsto\varphi\}\).
  \item If \(\mathfrak s=p_n\) or \(\mathfrak s=\bottom\), return the empty
        function if \(\varphi=\mathfrak s\), and \(\mathsf{fail}\) otherwise.
  \item If \(\mathfrak s=\neg\mathfrak s_1\): if \(\varphi\) is a negation
        \(\neg\varphi_1\), return \(\operatorname{match}(\mathfrak
        s_1,\varphi_1)\); otherwise return \(\mathsf{fail}\).
  \item If \(\mathfrak s=(\mathfrak s_1\circ\mathfrak s_2)\): if \(\varphi\)
        is a binary formula \((\varphi_1\circ'\varphi_2)\) with
        \(\circ'=\circ\), compute \(m_1=\operatorname{match}(\mathfrak
        s_1,\varphi_1)\) and \(m_2=\operatorname{match}(\mathfrak
        s_2,\varphi_2)\); if either is \(\mathsf{fail}\), or if some letter in
        the domain of both receives different formulas from \(m_1\) and
        \(m_2\), return \(\mathsf{fail}\); otherwise return \(m_1\cup m_2\).
        In every other case return \(\mathsf{fail}\).
\end{enumerate}
\emph{Side conditions.} The decomposition of \(\varphi\) in steps 3 and 4 is
the unique one given by \cref{thm:prop-unique-readability}; the procedure is
not defined on strings outside \(\mathsf{Fm}_{\mathsf{CPL}}\). The consistency
test in step 4 is what enforces uniform instantiation of a repeated letter.
Each recursive call is on an immediate subschema, so the recursion
terminates after at most as many calls as \(\mathfrak s\) has symbol
occurrences.
\end{hbmethod}

\begin{hbtheorem}{thm:p01-w02-p01-004-matching-correct}{proved=annex:p01-w02-p01-004-matching-correct}{Correctness of matching}
For every \(\mathfrak s\in\mathrm{Sch}\) and
\(\varphi\in\mathsf{Fm}_{\mathsf{CPL}}\):
\begin{enumerate}
  \item if \(\operatorname{match}(\mathfrak s,\varphi)=m\ne\mathsf{fail}\),
        then the domain of \(m\) is \(\operatorname{MV}(\mathfrak s)\) and
        \(\mathfrak s[\iota]=\varphi\) for every instantiation \(\iota\)
        extending \(m\);
  \item if \(\varphi\in\operatorname{Inst}(\mathfrak s)\), then
        \(\operatorname{match}(\mathfrak s,\varphi)\ne\mathsf{fail}\).
\end{enumerate}
Consequently \(\varphi\in\operatorname{Inst}(\mathfrak s)\) if and only if
\(\operatorname{match}(\mathfrak s,\varphi)\ne\mathsf{fail}\), and in that
case \(\operatorname{match}(\mathfrak s,\varphi)\) is the restriction to
\(\operatorname{MV}(\mathfrak s)\) of every instantiation \(\iota\) with
\(\mathfrak s[\iota]=\varphi\).
\end{hbtheorem}

\begin{hbtheorem}{thm:p01-w02-p01-004-instance-count}{proved=annex:p01-w02-p01-004-instance-count}{Size of instance sets}
Let \(\mathfrak s\in\mathrm{Sch}\).
\begin{enumerate}
  \item If \(\operatorname{MV}(\mathfrak s)=\varnothing\), then
        \(\mathfrak s\in\mathsf{Fm}_{\mathsf{CPL}}\) and
        \(\operatorname{Inst}(\mathfrak s)=\{\mathfrak s\}\).
  \item If \(\operatorname{MV}(\mathfrak s)\ne\varnothing\), then
        \(\operatorname{Inst}(\mathfrak s)\) is countably infinite.
\end{enumerate}
\end{hbtheorem}

The relation to uniform substitution passes through a translation of the
schemas into formulas, which keeps the fixed variables and the schematic
letters apart.
Let \(\operatorname{tr}\) be the map on strings over \(A^{+}\) that replaces
each token \(p_n\) by \(p_{2n}\) and each token \(\xi_k\) by \(p_{2k+1}\), and
leaves every other token unchanged.

\begin{hbtheorem}{thm:p01-w02-p01-004-substitution-link}{proved=annex:p01-w02-p01-004-substitution-link}{Instances as uniform-substitution instances}
For every instantiation \(\iota\), let \(\upsilon_\iota\) be the uniform
substitution with \(\upsilon_\iota(p_{2n})=p_n\) and
\(\upsilon_\iota(p_{2k+1})=\iota(\xi_k)\) for all \(n,k\in\N\). Then for every
\(\mathfrak s\in\mathrm{Sch}\):
\begin{enumerate}
  \item \(\operatorname{tr}(\mathfrak s)\in\mathsf{Fm}_{\mathsf{CPL}}\),
        and \(\mathfrak s[\iota]\) is the formula
        \(\operatorname{tr}(\mathfrak s)[\upsilon_\iota]\);
  \item \(\operatorname{Inst}(\mathfrak s)\) is the set of formulas
        \(\operatorname{tr}(\mathfrak s)[\upsilon]\) for the uniform
        substitutions \(\upsilon\) with \(\upsilon(p_{2n})=p_n\) for all
        \(n\);
  \item if \(\entails\operatorname{tr}(\mathfrak s)\), then
        \(\entails\varphi\) for every
        \(\varphi\in\operatorname{Inst}(\mathfrak s)\).
\end{enumerate}
\end{hbtheorem}

Clause 3 is the familiar use of a schema: a valid pattern has only valid
instances. Clause 2 shows that a schema without fixed variables \(p_n\) has
exactly the uniform-substitution instances of the formula
\(\operatorname{tr}(\mathfrak s)\) as its instances. This agrees with the
inherited rule sheets, where a displayed rule is a rule schema whose
instances arise by uniformly substituting formulas for its formula
metavariables \cite[ch.~``Rule Sheets'', section ``Classical natural
deduction'']{FrozenFormalLogicSchemasCodes}.

\subsection{Quotation}
\label{sec:p01-w02-p01-004-quotation}

A quotation name is formed in the metalanguage by enclosing an inscription in
quotation marks; it denotes that inscription and nothing else. The
metalinguistic sentence ``the string \(\text{`}(p_0\to p_0)\text{'}\) has
five symbol occurrences'' mentions one formula. Quotation names are names of
the metalanguage. A quotation cannot be filled in: the quotation of the
display \((\varphi\to\varphi)\) names a string containing the Greek letter,
not a formula, which is why the handbook states claims about all formulas of
a given shape by the bound reading of
\cref{def:p01-w02-p01-004-display-reading} rather than by quoting a display.

\begin{hbwarning}{def:p01-w02-p01-004-corner-quotes}{Corner quotes in the arithmetic chapters name codes}
\Cref{sec:levels-02} allows quotation marks or corner quotes for literal
mention. From \cref{sec:arithmetic-coding-04} onward, corner quotes have a
fixed different meaning: \(\ulcorner E\urcorner\) is the numerical code of the
expression \(E\) in its declared syntactic sort, a natural number, and
\(\overline{\ulcorner E\urcorner}\) is the arithmetic numeral of that number,
an object-language term. In those parts a corner-quoted expression is
therefore never a quotation name of an inscription. The code of this section
is written \(\#\varphi\) to keep it apart from the arithmetic coding.
\end{hbwarning}

\subsection{Numerical codes}
\label{sec:p01-w02-p01-004-codes}

A code replaces a formula by a natural number from which the formula can be
recovered. The construction below assigns a numerical tag to each formation
clause and combines the tag with the codes of the immediate subformulas by
the pairing of \cref{sec:strings-01}.

\begin{hbdefinition}{def:p01-w02-p01-004-code}{Numerical code of a propositional formula}
Write \(\pi\) for \(\operatorname{pair}_{\N}\). The \emph{code}
\(\#\varphi\in\N\) of a formula \(\varphi\in\mathsf{Fm}_{\mathsf{CPL}}\) is
defined by structural recursion:
\[
\begin{aligned}
 \#p_n&=1+\pi(0,n),&
 \#(\varphi\land\psi)&=1+\pi(3,\pi(\#\varphi,\#\psi)),\\
 \#\bottom&=1+\pi(1,0),&
 \#(\varphi\lor\psi)&=1+\pi(4,\pi(\#\varphi,\#\psi)),\\
 \#\neg\varphi&=1+\pi(2,\#\varphi),&
 \#(\varphi\to\psi)&=1+\pi(5,\pi(\#\varphi,\#\psi)).
\end{aligned}
\]
The first argument of the outer pairing is the \emph{tag} of the formation
clause. \Cref{thm:prop-unique-readability} makes the recursion well defined:
each formula falls under exactly one line, with unique immediate
subformulas.
\end{hbdefinition}

\begin{hbtheorem}{thm:p01-w02-p01-004-code-properties}{proved=annex:p01-w02-p01-004-code-properties}{Properties of the code}
\begin{enumerate}
  \item For all \(a,b\in\N\), \(\pi(a,b)\ge a\) and \(\pi(a,b)\ge b\).
  \item \(\#\varphi\ge1\) for every formula \(\varphi\).
  \item The map \(\#:\mathsf{Fm}_{\mathsf{CPL}}\to\N\) is injective.
  \item If \(\psi\) is an immediate subformula of \(\varphi\), then
        \(\#\psi<\#\varphi\).
\end{enumerate}
\end{hbtheorem}

\begin{hbmethod}{def:p01-w02-p01-004-decoding}{Recognizing and decoding codes}
\emph{Input:} \(c\in\N\). \emph{Output:} a formula \(\operatorname{dec}(c)\),
or the answer \(\mathsf{none}\). For \(z\in\N\), let \(\pi^{-1}(z)\) be the
unique pair \((a,b)\) with \(\pi(a,b)=z\), computed as in
\cref{sec:strings-01}: with \(t_s=s(s+1)/2\), find the unique \(s\) with
\(t_s\le z<t_{s+1}\), and put \(b=z-t_s\) and \(a=s-b\).
\begin{enumerate}
  \item If \(c=0\), return \(\mathsf{none}\). Otherwise let
        \((k,x)=\pi^{-1}(c-1)\).
  \item If \(k=0\), return \(p_x\).
  \item If \(k=1\), return \(\bottom\) when \(x=0\) and \(\mathsf{none}\)
        otherwise.
  \item If \(k=2\), return \(\neg\operatorname{dec}(x)\) when
        \(\operatorname{dec}(x)\) is a formula and \(\mathsf{none}\)
        otherwise.
  \item If \(k\in\{3,4,5\}\), let \((a,b)=\pi^{-1}(x)\); when
        \(\operatorname{dec}(a)\) and \(\operatorname{dec}(b)\) are both
        formulas return \((\operatorname{dec}(a)\circ_k\operatorname{dec}(b))\),
        where \(\circ_3={\land}\), \(\circ_4={\lor}\) and \(\circ_5={\to}\);
        otherwise return \(\mathsf{none}\).
  \item If \(k\ge6\), return \(\mathsf{none}\).
\end{enumerate}
\emph{Side conditions and termination.} The search for \(s\) in
\(\pi^{-1}(z)\) is bounded: \(t_s\ge s\) for every \(s\), so \(s\le z\). Every
recursive call is made at \(x\), \(a\) or \(b\), and each of these is at most
\(c-1\) by \cref{thm:p01-w02-p01-004-code-properties}, clause 1, so the
recursion is a course-of-values recursion on \(c\) and terminates on every
input.
\end{hbmethod}

\begin{hbtheorem}{thm:p01-w02-p01-004-decoding-correct}{proved=annex:p01-w02-p01-004-decoding-correct}{Correctness of decoding}
For every formula \(\varphi\), the decoding of its code is \(\varphi\), that
is, \(\operatorname{dec}(\#\varphi)=\varphi\). For every \(c\in\N\), if \(\operatorname{dec}(c)\) is a formula \(\varphi\), then
\(c=\#\varphi\). Consequently a natural number \(c\) is the code of a formula
if and only if \(\operatorname{dec}(c)\ne\mathsf{none}\), and
\(\operatorname{dec}\) restricted to the set of codes is the inverse of
\(\#\).
\end{hbtheorem}

\begin{hbexample}{def:p01-w02-p01-004-worked}{Coding, decoding and matching one formula}
Let \(\varphi=(p_0\to\neg p_0)\). Then
\[
\begin{aligned}
 \#p_0&=1+\pi(0,0)=1+0=1,\\
 \#\neg p_0&=1+\pi(2,1)=1+(t_3+1)=1+7=8,\\
 \pi(\#p_0,\#\neg p_0)&=\pi(1,8)=t_9+8=45+8=53,\\
 \#\varphi&=1+\pi(5,53)=1+(t_{58}+53)=1+(1711+53)=1765.
\end{aligned}
\]
Decoding \(1765\): \(c-1=1764\), and \(t_{58}=1711\le1764<1770=t_{59}\), so
\(b=53\) and \(a=58-53=5\). The tag \(5\) calls for an implication, and
\(\pi^{-1}(53)=(1,8)\) because \(t_9=45\le53<55=t_{10}\). Next
\(\pi^{-1}(0)=(0,0)\) gives \(\operatorname{dec}(1)=p_0\), and
\(\pi^{-1}(7)=(2,1)\), from \(t_3=6\le7<10=t_4\), gives
\(\operatorname{dec}(8)=\neg p_0\). Hence
\(\operatorname{dec}(1765)=(p_0\to\neg p_0)\). The tag \(1\) yields a code
only with second component \(0\): \(\pi^{-1}(1)=(1,0)\), so
\(\operatorname{dec}(2)=\bottom\) and \(\#\bottom=2\); but
\(\pi^{-1}(4)=(1,1)\), from \(t_2=3\le4<6=t_3\), has tag \(1\) and
\(x=1\ne0\), so \(\operatorname{dec}(5)=\mathsf{none}\) and \(5\) is not the
code of any formula.

Now match \(\varphi\) against the schema \(\mathfrak s=(\xi_0\to(\xi_1\to
\xi_0))\), displayed as \(\varphi\to(\psi\to\varphi)\). Step 4 of
\cref{def:p01-w02-p01-004-matching} compares the main connectives, both
\(\to\), and calls \(\operatorname{match}(\xi_0,p_0)=\{\xi_0\mapsto p_0\}\)
and \(\operatorname{match}((\xi_1\to\xi_0),\neg p_0)\). The second call
fails, since \(\neg p_0\) is a negation and not an implication; so
\(\varphi\notin\operatorname{Inst}(\mathfrak s)\). The formula
\((p_0\to(\neg p_0\to p_0))\) does match, with the unique restriction
\(\{\xi_0\mapsto p_0,\ \xi_1\mapsto\neg p_0\}\). The formula
\((p_0\to(p_1\to p_2))\) does not match: the two calls for \(\xi_0\) return
\(p_0\) and \(p_2\), and the consistency test of step 4 fails.
\end{hbexample}

\subsection{The schema/code distinction}
\label{sec:p01-w02-p01-004-distinction}

\begin{hbdefinition}{def:p01-w02-p01-004-distinction}{Schema, quotation name, code and numeral}
Four kinds of expression are used to speak about formulas.
\begin{enumerate}
  \item A \emph{formula schema} \(\mathfrak s\) is a metalinguistic
        expression; it stands for the family
        \(\operatorname{Inst}(\mathfrak s)\) of object-language formulas,
        its instances, and is not itself a formula when it contains a
        schematic letter.
  \item A \emph{quotation name} is a metalinguistic name of exactly one
        inscription.
  \item A \emph{code} is a natural number; the number \(\#\varphi\) is
        correlated with exactly one formula, \(\varphi\), by an injective
        coding fixed in advance.
  \item In an object language containing numerals, the \emph{numeral} of a
        code is a closed object-language term denoting that number; through
        the fixed coding, the object language can thereby mention one
        particular formula.
\end{enumerate}
\end{hbdefinition}

\begin{hbtheorem}{thm:p01-w02-p01-004-distinction}{proved=annex:p01-w02-p01-004-distinction}{Schemas stand for families; codes name single formulas}
\begin{enumerate}
  \item For every formula \(\varphi\) there is exactly one \(c\in\N\) with
        \(\operatorname{dec}(c)=\varphi\), namely \(c=\#\varphi\); and every
        \(c\) with \(\operatorname{dec}(c)\ne\mathsf{none}\) is correlated
        with exactly one formula.
  \item A schema \(\mathfrak s\) is a formula if and only if
        \(\operatorname{Inst}(\mathfrak s)\) has exactly one member, and
        then \(\operatorname{Inst}(\mathfrak s)=\{\mathfrak s\}\).
  \item If \(\operatorname{MV}(\mathfrak s)\ne\varnothing\), there is a
        schema \(\mathfrak s'\ne\mathfrak s\) with
        \(\operatorname{Inst}(\mathfrak s')=\operatorname{Inst}(\mathfrak
        s)\). Thus a schema is not recoverable from its instance family,
        whereas a formula is recoverable from its code.
\end{enumerate}
\end{hbtheorem}

\begin{hbtheorem}{thm:p01-w02-p01-004-instance-codes}{proved=annex:p01-w02-p01-004-instance-codes}{The codes of an instance family}
For every schema \(\mathfrak s\),
\[
 \{\#\varphi:\varphi\in\operatorname{Inst}(\mathfrak s)\}
 =\{c\in\N:\operatorname{dec}(c)\ne\mathsf{none}\text{ and }
    \operatorname{match}(\mathfrak s,\operatorname{dec}(c))\ne\mathsf{fail}\}.
\]
Membership of a number in this set is therefore decided by running the
procedures of \cref{def:p01-w02-p01-004-decoding} and
\cref{def:p01-w02-p01-004-matching} in turn.
\end{hbtheorem}

\Cref{thm:p01-w02-p01-004-instance-codes} is the point at which the two
devices meet. A schema is a pattern in the metalanguage, but the codes of its
instances form one set of natural numbers. An object language that can speak
about numbers can therefore speak about the instance family only through such
a set of codes, while the schema itself remains outside it. Coding the
schema itself is possible, by extending the tags to the letters \(\xi_k\);
the resulting number names the one string \(\mathfrak s\), not its instances.

\begin{hbwarning}{def:p01-w02-p01-004-schema-not-sentence}{A schema is not one sentence, and a code does not supply one}
A schema with a schematic letter is a family of formulas, not a single
formula, and an axiom schema adds every instance as an axiom; the inherited
checklist for axiomatizations requires that a schema be identified as a set
of sentences \cite[ch.~``Theories and Axiomatization'', warning on writing an
axiomatization, item~6]{FrozenFormalLogicSchemasCodes}. Codes do not remove
this distinction. In the arithmetic language \(\Lang_A\) with the coding of
\cref{sec:arithmetic-coding-04}, consider the display
\[
 \mathsf{Tr}(\overline{\ulcorner\sigma\urcorner})\leftrightarrow\sigma .
\]
Here \(\sigma\) occurs twice in different roles: on the right it is a
metavariable in the schematic reading, whose instances are sentences used in
the formula; on the left it occurs inside a numeral, which is one closed
object-language term mentioning the instance sentence through its code. For
a fixed \(\Lang_A\)-formula \(\mathsf{Tr}(x)\), asking that every instance be
true in the structure \(\mathcal N\) of the natural numbers is the definition
of an arithmetical truth definition, and \cref{thm:tarski-undefinability}
shows that no \(\Lang_A\)-formula satisfies it. The family of truths named by
the schema cannot be captured by one object-language formula applied to
codes. The inherited remark on partial truth definitions records that this
obstacle concerns one formula for all arithmetic sentences at once, not one
formula for each fixed level of the arithmetical hierarchy
\cite[ch.~``Provability, the Second Incompleteness Theorem, and Truth'',
remark ``Partial truth definitions'']{FrozenFormalLogicSchemasCodes}.
\end{hbwarning}

\paragraph{Limitations.} The schemas of this section have propositional
instances only. Schemas with term and variable letters, as in the quantifier
rules of first-order natural deduction, need further side conditions on
instantiation (freshness of eigenvariables and freedom of a term for a
variable), so their instance sets are not obtained by plain instantiation as
in \cref{def:p01-w02-p01-004-instance}. The decoding procedure is shown to
terminate and to be correct; its classification as a primitive-recursive
function requires the class of primitive-recursive functions of
\cref{sec:turing-machines-02} and is carried out for the arithmetic coding in \cref{sec:arithmetic-coding-04}, not for the code
\(\#\) defined here.

\paragraph{Relationships.} The code \(\#\) and the arithmetic coding of
\cref{sec:arithmetic-coding-04} use the same pairing function and the same
pattern of a constructor tag combined with child codes, and both rely on
decrease of child codes for course-of-values recursion; the arithmetic coding
adds terms, quantifiers, symbols and strings as further sorts. The injection
of finite strings into \(\N\) in \cref{sec:strings-01} codes every string,
formula or not, and so does not by itself give the clause-by-clause recovery
of \cref{thm:p01-w02-p01-004-decoding-correct}. Instances of schemas are
uniform-substitution instances
(\cref{thm:p01-w02-p01-004-substitution-link}), so
\cref{thm:p02-uniform-substitution} transfers validity from a schema's
translation to all its instances. The level error recorded in the inherited
error taxonomy, the confusion of a formula with a metalinguistic statement,
is repaired by quotation or by quantifying with a metavariable
\cite[ch.~``Writing Formal Definitions and Proofs'', section ``Error
taxonomy'']{FrozenFormalLogicSchemasCodes}; the bound and schematic readings
of \cref{def:p01-w02-p01-004-display-reading} make the second repair exact.
